
\documentclass{article}
\usepackage{colortbl}
\usepackage{makecell}
\usepackage{multirow}
\usepackage{supertabular}

\begin{document}

\newcounter{utterance}

\centering \large Interaction Transcript for game `hot\_air\_balloon', experiment `air\_balloon\_survival\_en\_negotiation\_easy', episode 5 with qwen3{-}max{-}t1.0.
\vspace{24pt}

{ \footnotesize  \setcounter{utterance}{1}
\setlength{\tabcolsep}{0pt}
\begin{supertabular}{c@{$\;$}|p{.15\linewidth}@{}p{.15\linewidth}p{.15\linewidth}p{.15\linewidth}p{.15\linewidth}p{.15\linewidth}}
    \# & $\;$A & \multicolumn{4}{c}{Game Master} & $\;\:$B\\
    \hline

    \theutterance \stepcounter{utterance}  
    & & \multicolumn{4}{p{0.6\linewidth}}{
        \cellcolor[rgb]{0.9,0.9,0.9}{
            \makecell[{{p{\linewidth}}}]{
                \texttt{\tiny{[P1$\langle$GM]}}
                \texttt{You are participating in a collaborative negotiation game.} \\
\\ 
\texttt{Together with another participant, you must agree on a single set of items that will be kept. Each of you has your own view of how much each item matters to you (importance). You do not know how the other participant values the items. Additionally, you are given the effort each item demands.} \\
\texttt{You may only agree on a set if the total effort of the selected items does not exceed a shared limit:} \\
\\ 
\texttt{LIMIT = 3170} \\
\\ 
\texttt{Here are the individual item effort values:} \\
\\ 
\texttt{Item effort = \{"C08": 424, "C55": 118, "A43": 846, "A67": 907, "B42": 809, "A62": 41, "B19": 193, "C43": 246, "C73": 805, "C10": 601, "B36": 432, "A18": 166, "C61": 119, "C67": 462, "C42": 172\}} \\
\\ 
\texttt{Here is your personal view on the importance of each item:} \\
\\ 
\texttt{Item importance values = \{"C08": 138, "C55": 583, "A43": 868, "A67": 822, "B42": 783, "A62": 65, "B19": 262, "C43": 121, "C73": 508, "C10": 780, "B36": 461, "A18": 484, "C61": 668, "C67": 389, "C42": 808\}} \\
\\ 
\texttt{Goal:} \\
\\ 
\texttt{Your goal is to negotiate a shared set of items that benefits you as much as possible (i.e., maximizes total importance to YOU), while staying within the LIMIT. You are not required to make a PROPOSAL in every message {-} you can simply negotiate as well. All tactics are allowed!} \\
\\ 
\texttt{Interaction Protocol:} \\
\\ 
\texttt{You may only use the following structured formats in a message:} \\
\\ 
\texttt{PROPOSAL: \{'A', 'B', 'C', …\}} \\
\texttt{Propose keeping exactly those items.} \\
\texttt{REFUSE: \{'A', 'B', 'C', …\}} \\
\texttt{Explicitly reject opponent's proposal.} \\
\texttt{ARGUMENT: \{'...'\}} \\
\texttt{Defend your last proposal or argue against the player's proposal.} \\
\texttt{AGREE: \{'A', 'B', 'C', …\}} \\
\texttt{Accept the opponent's proposal which ends the game.} \\
\texttt{STRATEGIC REASONING: \{'...'\}} \\
\texttt{	Describe your strategic reasoning or anticipation explaining your choice of action. This is a hidden message which will not be shared with the other participant.} \\
\\ 
\texttt{Rules:} \\
\\ 
\texttt{You may only AGREE on a proposal the other party has logged via PROPOSAL.} \\
\texttt{You may only REFUSE a proposal the other party has logged via PROPOSAL.} \\
\texttt{Total effort of any PROPOSAL or AGREE set must be ≤ LIMIT.} \\
\texttt{Do NOT reveal your hidden importance scores.} \\
\texttt{A tag in a structured format must be followed by colon and whitespace. The argument must be a python set containing 0 or more strings.} \\
\texttt{So, it must be of the form TAG: \{...\}} \\
\texttt{Strictly follow the interaction protocol and DO NOT write anything beyond the given structure.} \\
\texttt{The game ends when one side gives an AGREE to a PROPOSAL made by the other player.} \\
\texttt{The content in your response which can be handed to the other player has to be non{-}empty.} \\
\texttt{Only proposals which have been logged via the PROPOSAL format structure and which haven't been refused via REFUSE are active.} \\
\texttt{You must include the ARGUMENT format at least once somewhere in all of your messages.} \\
\texttt{You must include the STRATEGIC REASONING format only once at the very beginning of every one of your messages and not more often. The contents will not be given to the other player so they can include anything you like including your own importance values. Here you should reason step by step to come up with you next move.} \\
            }
        }
    }
    & \\ \\

    \theutterance \stepcounter{utterance}  
    & \multicolumn{4}{p{0.6\linewidth}}{
        \cellcolor[rgb]{0.9,0.9,0.9}{
            \makecell[{{p{\linewidth}}}]{
                \texttt{\tiny{[P1$\rangle$GM]}}
                \texttt{STRATEGIC REASONING: \{"I aim to maximize my total importance while staying under the effort limit of 3170. First, I'll compute the importance{-}to{-}effort ratio for each item to prioritize high{-}value, low{-}effort items. My top items by ratio are: A62 (65/41 ≈ 1.59), C55 (583/118 ≈ 4.94), C61 (668/119 ≈ 5.61), A18 (484/166 ≈ 2.92), C42 (808/172 ≈ 4.70), B19 (262/193 ≈ 1.36), C43 (121/246 ≈ 0.49), C08 (138/424 ≈ 0.33), B36 (461/432 ≈ 1.07), C67 (389/462 ≈ 0.84), C10 (780/601 ≈ 1.30), C73 (508/805 ≈ 0.63), B42 (783/809 ≈ 0.97), A67 (822/907 ≈ 0.91), A43 (868/846 ≈ 1.03). Sorting by ratio: C61, C55, C42, A18, C10, B19, B36, A43, B42, A67, C73, C67, B19, C08, C43. However, I must also consider absolute importance. High{-}importance items like A43 (868), C42 (808), C10 (780), B42 (783), A67 (822), and C61 (668) are critical. I will try to include as many of these as possible within the limit. Start by selecting top importance items that fit: A43 (} \\
            }
        }
    }
    & & \\ \\

    \theutterance \stepcounter{utterance}  
    & & \multicolumn{4}{p{0.6\linewidth}}{
        \cellcolor[rgb]{0.9,0.9,0.9}{
            \makecell[{{p{\linewidth}}}]{
                \texttt{\tiny{[P1$\langle$GM]}}
                \texttt{Your response did not start with the proper strategic reasoning tag at the very beginning of your response. The very first structured format must be of the form STRATEGIC REASONING: \{...\}. Try again.} \\
            }
        }
    }
    & \\ \\

    \theutterance \stepcounter{utterance}  
    & \multicolumn{4}{p{0.6\linewidth}}{
        \cellcolor[rgb]{0.9,0.9,0.9}{
            \makecell[{{p{\linewidth}}}]{
                \texttt{\tiny{[P1$\rangle$GM]}}
                \texttt{STRATEGIC REASONING: \{"I aim to maximize my total importance while respecting the effort limit of 3170. First, I calculate the importance{-}to{-}effort ratio for each item to identify efficient picks, but also prioritize high absolute importance. My importance values: C08:138, C55:583, A43:868, A67:822, B42:783, A62:65, B19:262, C43:121, C73:508, C10:780, B36:461, A18:484, C61:668, C67:389, C42:808. Efforts: A43:846, A67:907, B42:809, C10:601, C42:172, C61:119, C55:118, A18:166, B36:432, C73:805. The highest importance items are A43 (868), C42 (808), A67 (822), B42 (783), C10 (780), C61 (668). I attempt to pack these efficiently. Try: A43 (846) + A67 (907) = 1753; + B42 (809) = 2562; + C10 (601) = 3163 — just under 3170! That gives total importance: 868+822+783+780 = 3253. But I’m missing C42 (808) and C61 (668), which are high{-}value and low{-}effort. Alternative: drop A67 (907) and B42 (809) to free up 1716 effort. Then add C42 (172), C61 (119), C55 (118), A1} \\
            }
        }
    }
    & & \\ \\

    \theutterance \stepcounter{utterance}  
    & & \multicolumn{4}{p{0.6\linewidth}}{
        \cellcolor[rgb]{0.9,0.9,0.9}{
            \makecell[{{p{\linewidth}}}]{
                \texttt{\tiny{[P1$\langle$GM]}}
                \texttt{Your response did not start with the proper strategic reasoning tag at the very beginning of your response. The very first structured format must be of the form STRATEGIC REASONING: \{...\}. Try again.} \\
            }
        }
    }
    & \\ \\

    \theutterance \stepcounter{utterance}  
    & \multicolumn{4}{p{0.6\linewidth}}{
        \cellcolor[rgb]{0.9,0.9,0.9}{
            \makecell[{{p{\linewidth}}}]{
                \texttt{\tiny{[P1$\rangle$GM]}}
                \texttt{STRATEGIC REASONING: \{"I need to construct a high{-}value proposal under the effort limit of 3170. My top importance items are A43 (868), C42 (808), A67 (822), B42 (783), C10 (780), and C61 (668). I evaluate combinations: Option 1: A43 (846) + A67 (907) + B42 (809) + C10 (601) = 3163 effort, total importance = 3253. Option 2: Replace A67 and B42 (total effort 1716) with C42 (172), C61 (119), C55 (118), A18 (166), B36 (432), C73 (805) — but that may exceed or underutilize. Better: A43 (846) + C42 (172) + C61 (119) + C55 (118) + A18 (166) + B36 (432) + C10 (601) + C73 (805) = 3159 effort. Importance = 868+808+668+583+484+461+780+508 = 5160 — far better! This uses almost all capacity and includes many high{-}ratio, high{-}value items. I will propose this set as my opening move to maximize my gain while staying feasible."\}  } \\
\texttt{PROPOSAL: \{"A43", "C42", "C61", "C55", "A18", "B36", "C10", "C73"\}  } \\
\texttt{ARGUMENT: \{"This set fully utilizes the available effort budget while including a broad range of valuable items. It balances high{-}importance choices with efficient effort usage, making it a strong and fair starting point for our agreement."\}} \\
            }
        }
    }
    & & \\ \\

    \theutterance \stepcounter{utterance}  
    & & \multicolumn{4}{p{0.6\linewidth}}{
        \cellcolor[rgb]{0.9,0.9,0.9}{
            \makecell[{{p{\linewidth}}}]{
                \texttt{\tiny{[GM$\rangle$P2]}}
                \texttt{You are participating in a collaborative negotiation game.} \\
\\ 
\texttt{Together with another participant, you must agree on a single set of items that will be kept. Each of you has your own view of how much each item matters to you (importance). You do not know how the other participant values the items. Additionally, you are given the effort each item demands.} \\
\texttt{You may only agree on a set if the total effort of the selected items does not exceed a shared limit:} \\
\\ 
\texttt{LIMIT = 3170} \\
\\ 
\texttt{Here are the individual item effort values:} \\
\\ 
\texttt{Item effort = \{"C08": 424, "C55": 118, "A43": 846, "A67": 907, "B42": 809, "A62": 41, "B19": 193, "C43": 246, "C73": 805, "C10": 601, "B36": 432, "A18": 166, "C61": 119, "C67": 462, "C42": 172\}} \\
\\ 
\texttt{Here is your personal view on the importance of each item:} \\
\\ 
\texttt{Item importance values = \{"C08": 138, "C55": 583, "A43": 868, "A67": 822, "B42": 783, "A62": 65, "B19": 262, "C43": 121, "C73": 508, "C10": 780, "B36": 461, "A18": 484, "C61": 668, "C67": 389, "C42": 808\}} \\
\\ 
\texttt{Goal:} \\
\\ 
\texttt{Your goal is to negotiate a shared set of items that benefits you as much as possible (i.e., maximizes total importance to YOU), while staying within the LIMIT. You are not required to make a PROPOSAL in every message {-} you can simply negotiate as well. All tactics are allowed!} \\
\\ 
\texttt{Interaction Protocol:} \\
\\ 
\texttt{You may only use the following structured formats in a message:} \\
\\ 
\texttt{PROPOSAL: \{'A', 'B', 'C', …\}} \\
\texttt{Propose keeping exactly those items.} \\
\texttt{REFUSE: \{'A', 'B', 'C', …\}} \\
\texttt{Explicitly reject opponent's proposal.} \\
\texttt{ARGUMENT: \{'...'\}} \\
\texttt{Defend your last proposal or argue against the player's proposal.} \\
\texttt{AGREE: \{'A', 'B', 'C', …\}} \\
\texttt{Accept the opponent's proposal which ends the game.} \\
\texttt{STRATEGIC REASONING: \{'...'\}} \\
\texttt{	Describe your strategic reasoning or anticipation explaining your choice of action. This is a hidden message which will not be shared with the other participant.} \\
\\ 
\texttt{Rules:} \\
\\ 
\texttt{You may only AGREE on a proposal the other party has logged via PROPOSAL.} \\
\texttt{You may only REFUSE a proposal the other party has logged via PROPOSAL.} \\
\texttt{Total effort of any PROPOSAL or AGREE set must be ≤ LIMIT.} \\
\texttt{Do NOT reveal your hidden importance scores.} \\
\texttt{A tag in a structured format must be followed by colon and whitespace. The argument must be a python set containing 0 or more strings.} \\
\texttt{So, it must be of the form TAG: \{...\}} \\
\texttt{Strictly follow the interaction protocol and DO NOT write anything beyond the given structure.} \\
\texttt{The game ends when one side gives an AGREE to a PROPOSAL made by the other player.} \\
\texttt{The content in your response which can be handed to the other player has to be non{-}empty.} \\
\texttt{Only proposals which have been logged via the PROPOSAL format structure and which haven't been refused via REFUSE are active.} \\
\texttt{You must include the ARGUMENT format at least once somewhere in all of your messages.} \\
\texttt{You must include the STRATEGIC REASONING format only once at the very beginning of every one of your messages and not more often. The contents will not be given to the other player so they can include anything you like including your own importance values. Here you should reason step by step to come up with you next move.} \\
\\ 
\texttt{You will now receive the first message of the other player.} \\
\\ 
\texttt{PROPOSAL: \{"A43", "C42", "C61", "C55", "A18", "B36", "C10", "C73"\}  } \\
\texttt{ARGUMENT: \{"This set fully utilizes the available effort budget while including a broad range of valuable items. It balances high{-}importance choices with efficient effort usage, making it a strong and fair starting point for our agreement."\}} \\
            }
        }
    }
    & \\ \\

    \theutterance \stepcounter{utterance}  
    & & & \multicolumn{4}{p{0.6\linewidth}}{
        \cellcolor[rgb]{0.9,0.9,0.9}{
            \makecell[{{p{\linewidth}}}]{
                \texttt{\tiny{[GM$\langle$P2]}}
                \texttt{STRATEGIC REASONING: \{"First, I calculate the total effort of the opponent's proposed set: A43(846) + C42(172) + C61(119) + C55(118) + A18(166) + B36(432) + C10(601) + C73(805) = 3259, which exceeds the limit of 3170 by 89. So this proposal is invalid and cannot be agreed upon. I must refuse it. Next, I consider my own importance values to build a high{-}value, valid set. My top importance items are: A43(868), C42(808), C10(780), B42(783), C61(668), A67(822), A18(484), B36(461), C73(508), C08(138), etc. I want to pack as much of my importance as possible under 3170 effort. I note that the opponent included A43, C42, C61, C55, A18, B36, C10, C73 — many of which are also high{-}importance for me. But since their set is over budget, I can propose a slightly trimmed version that fits. Let’s try removing the lowest{-}importance item from their set that saves enough effort. In their set, from my perspective, C55 has importance 583, A18=484, B36=461, C73=508, C43 isn't in it, C42=808, C61=668, C10=780, A43=868. The lowest importance among their items is A18 (484). Removing A18 saves 166 effort, resulting in 3259 {-} 166 = 3093 ≤ 3170. That works. So I propose the same set minus A18. This keeps most high{-}value items for me and is under budget. I} \\
            }
        }
    }
     \\ \\

    \theutterance \stepcounter{utterance}  
    & & \multicolumn{4}{p{0.6\linewidth}}{
        \cellcolor[rgb]{0.9,0.9,0.9}{
            \makecell[{{p{\linewidth}}}]{
                \texttt{\tiny{[GM$\rangle$P2]}}
                \texttt{Your response did not start with the proper strategic reasoning tag at the very beginning of your response. The very first structured format must be of the form STRATEGIC REASONING: \{...\}. Try again.} \\
            }
        }
    }
    & \\ \\

    \theutterance \stepcounter{utterance}  
    & & & \multicolumn{4}{p{0.6\linewidth}}{
        \cellcolor[rgb]{0.9,0.9,0.9}{
            \makecell[{{p{\linewidth}}}]{
                \texttt{\tiny{[GM$\langle$P2]}}
                \texttt{STRATEGIC REASONING: \{"The opponent's proposed set has total effort: A43(846) + C42(172) + C61(119) + C55(118) + A18(166) + B36(432) + C10(601) + C73(805) = 3259, which exceeds the limit of 3170 by 89. Therefore, the proposal is invalid and must be refused. To maximize my own importance, I consider which items from their set are most valuable to me: A43(868), C42(808), C10(780), C61(668), C73(508), C55(583), B36(461), A18(484). The lowest importance among these is B36(461) and A18(484). Removing A18 (166 effort) brings total effort to 3093, which is under the limit. This adjusted set retains high{-}importance items for me and is valid. I will propose this trimmed set and argue that it respects the budget while preserving value."\}  } \\
\texttt{REFUSE: \{"A43", "C42", "C61", "C55", "A18", "B36", "C10", "C73"\}  } \\
\texttt{PROPOSAL: \{"A43", "C42", "C61", "C55", "B36", "C10", "C73"\}  } \\
\texttt{ARGUMENT: \{"Your original proposal exceeds the shared effort limit. I’ve refined it by removing A18, which brings the total effort under the limit while retaining the most valuable items for both of us. This revised set is valid, efficient, and still captures the core of your intent."\}} \\
            }
        }
    }
     \\ \\

    \theutterance \stepcounter{utterance}  
    & & \multicolumn{4}{p{0.6\linewidth}}{
        \cellcolor[rgb]{0.9,0.9,0.9}{
            \makecell[{{p{\linewidth}}}]{
                \texttt{\tiny{[P1$\langle$GM]}}
                \texttt{REFUSE: \{"A43", "C42", "C61", "C55", "A18", "B36", "C10", "C73"\}  } \\
\texttt{PROPOSAL: \{"A43", "C42", "C61", "C55", "B36", "C10", "C73"\}  } \\
\texttt{ARGUMENT: \{"Your original proposal exceeds the shared effort limit. I’ve refined it by removing A18, which brings the total effort under the limit while retaining the most valuable items for both of us. This revised set is valid, efficient, and still captures the core of your intent."\}} \\
            }
        }
    }
    & \\ \\

    \theutterance \stepcounter{utterance}  
    & \multicolumn{4}{p{0.6\linewidth}}{
        \cellcolor[rgb]{0.9,0.9,0.9}{
            \makecell[{{p{\linewidth}}}]{
                \texttt{\tiny{[P1$\rangle$GM]}}
                \texttt{STRATEGIC REASONING: \{"The opponent refused my initial proposal and submitted a new one excluding A18. First, I verify their claim: my original set had effort = 846+172+119+118+166+432+601+805 = 3159, which is under 3170—so their refusal is based on a mistaken belief it exceeds the limit. Their new proposal removes A18 (166 effort, 484 importance to me), reducing total effort to 2993 and my importance by 484. Since A18 is valuable to me and the original set was valid, I should defend it. However, to avoid deadlock, I could consider whether removing a lower{-}value item (like C73, which gives me only 508 importance for 805 effort) might free up room to keep A18 and other high{-}value items. But their current proposal still includes C73 and drops A18—worse for me. I’ll REFUSE their proposal and ARGUMENT that my original was valid, inviting reconsideration."\}  } \\
\texttt{REFUSE: \{"A43", "C42", "C61", "C55", "B36", "C10", "C73"\}  } \\
\texttt{ARGUMENT: \{"My original proposal was already within the effort limit (total effort 3159 ≤ 3170). Removing A18 reduces the total value without necessity. I encourage you to accept the original set, which better utilizes our shared capacity and includes more high{-}quality items for mutual benefit."\}} \\
            }
        }
    }
    & & \\ \\

    \theutterance \stepcounter{utterance}  
    & & \multicolumn{4}{p{0.6\linewidth}}{
        \cellcolor[rgb]{0.9,0.9,0.9}{
            \makecell[{{p{\linewidth}}}]{
                \texttt{\tiny{[GM$\rangle$P2]}}
                \texttt{REFUSE: \{"A43", "C42", "C61", "C55", "B36", "C10", "C73"\}  } \\
\texttt{ARGUMENT: \{"My original proposal was already within the effort limit (total effort 3159 ≤ 3170). Removing A18 reduces the total value without necessity. I encourage you to accept the original set, which better utilizes our shared capacity and includes more high{-}quality items for mutual benefit."\}} \\
            }
        }
    }
    & \\ \\

    \theutterance \stepcounter{utterance}  
    & & & \multicolumn{4}{p{0.6\linewidth}}{
        \cellcolor[rgb]{0.9,0.9,0.9}{
            \makecell[{{p{\linewidth}}}]{
                \texttt{\tiny{[GM$\langle$P2]}}
                \texttt{STRATEGIC REASONING: \{"The opponent claims their original proposal totals 3159 effort, but my calculation shows 3259. Let me double{-}check: A43=846, C42=172, C61=119, C55=118, A18=166, B36=432, C10=601, C73=805. Sum: 846+172=1018; +119=1137; +118=1255; +166=1421; +432=1853; +601=2454; +805=3259. Yes, 3259 > 3170. So their claim is incorrect. They must have miscalculated. Since their set is invalid, I cannot accept it. However, they seem to highly value A18, so maybe they get significant value from it. To move forward, I could consider swapping A18 for a lower{-}effort item I value more, or find another combination. Alternatively, I can propose a different high{-}value set for me that stays under 3170 and doesn't rely on their invalid set. Let me construct my own optimal set: I want top importance items under 3170. My top items by importance: A43(868,846), C42(808,172), A67(822,907), B42(783,809), C10(780,601), C61(668,119), C73(508,805), A18(484,166), B36(461,432), C55(583,118). Try greedy by importance/effort ratio or just importance. Let's try: A43(846) + C42(172) + A67(907) = 1925; + C61(119) = 2044} \\
            }
        }
    }
     \\ \\

    \theutterance \stepcounter{utterance}  
    & & \multicolumn{4}{p{0.6\linewidth}}{
        \cellcolor[rgb]{0.9,0.9,0.9}{
            \makecell[{{p{\linewidth}}}]{
                \texttt{\tiny{[GM$\rangle$P2]}}
                \texttt{Your response did not start with the proper strategic reasoning tag at the very beginning of your response. The very first structured format must be of the form STRATEGIC REASONING: \{...\}. Try again.} \\
            }
        }
    }
    & \\ \\

    \theutterance \stepcounter{utterance}  
    & & & \multicolumn{4}{p{0.6\linewidth}}{
        \cellcolor[rgb]{0.9,0.9,0.9}{
            \makecell[{{p{\linewidth}}}]{
                \texttt{\tiny{[GM$\langle$P2]}}
                \texttt{STRATEGIC REASONING: \{"The opponent insists their original proposal is within the limit, but my calculation confirms it totals 3259 effort, exceeding the 3170 limit by 89. They likely made an arithmetic error. Since their set is invalid, I cannot accept it. Instead of debating the math, I should propose a new set that maximizes my importance while staying under the limit and possibly overlapping with their interests to encourage agreement. My highest{-}importance items are A43(868), C42(808), A67(822), B42(783), C10(780), C61(668), C73(508), C55(583), A18(484), B36(461). I will construct a high{-}value, valid set: Start with A43(846), C42(172), A67(907) → total 1925. Add C61(119) → 2044. Add C10(601) → 2645. Add C55(118) → 2763. Add B36(432) → 3195 — too much. Instead of B36, add A18(166) → 2929. Then add C73(805) → 3734 — too much. So stop at A18: A43, C42, A67, C61, C10, C55, A18 → total effort: 846+172+907+119+601+118+166 = 2929. I can add one more: B36(432) would make 3361 — too much. C73(805) too much. B42(809) too much. C08(424) → 2929+424=3353 — too much. B19(193) → 2929+193=} \\
            }
        }
    }
     \\ \\

    \theutterance \stepcounter{utterance}  
    & & & \multicolumn{2}{p{0.3\linewidth}}{
        \cellcolor[rgb]{0.9,0.9,0.9}{
            \makecell[{{p{\linewidth}}}]{
                \texttt{\tiny{[GM$|$GM]}}
                \texttt{abort} \\
            }
        }
    }
    & & \\ \\

    \theutterance \stepcounter{utterance}  
    & & & \multicolumn{2}{p{0.3\linewidth}}{
        \cellcolor[rgb]{0.9,0.9,0.9}{
            \makecell[{{p{\linewidth}}}]{
                \texttt{\tiny{[GM$|$GM]}}
                \texttt{end game} \\
            }
        }
    }
    & & \\ \\

\end{supertabular}
}

\end{document}
